\documentclass[10pt,a4paper]{amsart}
\usepackage[margin=19mm]{geometry}
\usepackage{amsmath,amssymb,mathtools}
\usepackage[scr=rsfs,cal=euler]{mathalfa}
\usepackage{xfrac}
\usepackage[unicode,hidelinks,bookmarks=false]{hyperref}
\newcommand{\ie}{i.e.\ }
\newcommand{\cf}{cf.\ }
\newcommand{\resp}{respectively\ }
\newcommand{\Subsection}{\subsection}
\providecommand{\nobreakdash}{\nobreak\hbox{-}}
\setlength{\emergencystretch}{2em}
\multlinegap0pt
\usepackage{amssymb}
\usepackage{xcolor}
\usepackage[shortlabels]{enumitem}
\usepackage{float}
\usepackage{tikz}
\usepackage{algorithm}\usepackage{algpseudocode}
\newtheorem{theorem}{Theorem}
\newtheorem{corollary}{Corollary}
\newtheorem{lemma}{Lemma}
\usepackage{cleveref}
\crefname{theorem}{Theorem}{Theorems}
\crefname{corollary}{Corollary}{Corollarys}
\crefname{lemma}{Lemma}{Lemmas}
\DeclareMathOperator*{\argmin}{argmin}
\newcommand{\bregaux}{r}
\newcommand{\bregprox}{g}
\newcommand{\bregsm}{f}
\newcommand{\cC}{{\mathcal{C}}}
\newcommand{\compidx}{%
  \mathord{\mathchoice
    {\vcenter{\hbox{\rotatebox{45}{\rule{0.39em}{0.39em}}}}}
    {\vcenter{\hbox{\rotatebox{45}{\rule{0.39em}{0.39em}}}}}
    {\vcenter{\hbox{\rotatebox{45}{\rule{0.29em}{0.29em}}}}}
    {\vcenter{\hbox{\rotatebox{45}{\rule{0.23em}{0.23em}}}}}}%
}
\newcommand{\inprod}[2]{\left\langle #1,#2 \right\rangle}
\newcommand{\pr}[1]{\left( #1 \right)}
\newcommand{\reals}{\mathbb{R}}
\newcommand{\tnabla}{\widetilde{\nabla}}
\title{Toward a Systematic Understanding and Interactive Search of \\ Lyapunov-Style Proofs in Optimization}
\author{TaeHo Yoon, Jaewook J. Suh, Edward Duc Hien Nguyen, Bicheng Ying, Shiqian Ma}
\date{Source: arXiv:2606.26077v1 (CC BY 4.0)}
\begin{document}
\maketitle

\noindent\emph{Source and licence.} Toward a Systematic Understanding and Interactive Search of \\ Lyapunov-Style Proofs in Optimization. Version arXiv:2606.26077v1 (CC BY 4.0), distributed by arXiv under the Creative Commons Attribution 4.0 International licence, http://creativecommons.org/licenses/by/4.0/, as recorded in the arXiv licence metadata for this exact version and shown on the abstract page https://arxiv.org/abs/2606.26077v1. Full licence notice: \texttt{licenses/arxiv-2606.26077-CC-BY-4.0.txt}.\par\smallskip

\noindent\textit{Editorial scope.} Counted unit: the article's lemma lemma:bppm-tightness (source0.tex lines 2940--2950). The article's complete original proof of it is delivered with it (source0.tex lines 3563--3742, one \texttt{proof} environment of 180 lines, which the article places away from the statement). No mathematics is written, repaired or re-derived: the statement and the proof are the article's own lines, in the article's own notation, and no retained line is substituted or re-flowed.  Retained uncounted as the source-local context the proof needs: lines 2706--2714; lines 2716--2720; lines 2805--2819; lines 2836--2840; lines 2843--2855; lines 2906--2927.  Nothing was omitted from any retained span, so the package carries no omission record.  No retained line contains a citation, so the wrapper carries no bibliography and no bibliography entry.  No retained line contains a figure environment, an image-inclusion command or drawing code, so no image is delivered and the package declares no asset dependency.  Editorial formatting only, in addition: an AMS \texttt{amsart} preamble that restates the article's own declarations and macros used by the retained lines, namely the environments \texttt{theorem}, \texttt{corollary}, \texttt{lemma} on separate counters, together with the standard packages those lines need.  The retained environments print in this document as Theorem 1, Corollary 1, Lemma 1 in order of appearance, whereas the article prints the same items with the numbering of its own full document.  The complete native source of the article as distributed in the arXiv e-print for this exact version is bundled unchanged as \texttt{sources/203666/source0.tex}.
\begin{align}
\label{eq:bpgm_main}
\tag{BPGM}
    x_{k+1}\in \argmin_{y\in C}
    \left\{
        \bregprox(y)+\inprod{\nabla \bregsm(x_k)}{y-x_k}
        +L D_h(y,x_k)
    \right\} 
\end{align}

\begin{align}
\label{eq:bpgm_optimality_main}
    \tnabla \bregprox(x_{k+1})+\nabla \bregsm(x_k)
    =-L\pr{\nabla h(x_{k+1})-\nabla h(x_k)}.
\end{align}

\begin{align}
\label{eq:bpgm_difference_identity}
\begin{aligned}
    \widetilde V_k-\widetilde V_{k+1}
    & =
    -\frac{k}{N}\cC_{\bregsm}(x_k,x_{k+1})
    -\frac{1}{N}\cC_{\bregsm}(x_{\compidx},x_k) \\
    & \quad 
    -\frac{k}{N}\cC_{\bregprox}(x_k,x_{k+1})
    -\frac{1}{N}\cC_{\bregprox}(x_{\compidx},x_{k+1}) \\
    & \quad 
    -\frac{Lk}{N}\cC_{\bregaux}(x_k,x_{k+1})
    -\frac{L(k+1)}{N} \cC_{\bregaux}(x_{k+1},x_k)
\end{aligned}
\end{align}

\begin{align}
\label{eqn:bpgm-lyapunov}
    \widetilde V_k=\frac{k}{N}\pr{F(x_k)-F(x_{\compidx})}+\frac{L}{N}D_h(x_{\compidx},x_k),
    \qquad k=0,\dots,N.
\end{align}

\begin{theorem}
\label{theorem:bpgm-lyapunov}
Let $h$ be a Legendre-type convex function with $C=\operatorname{int}\operatorname{dom}h\neq\emptyset$.
Let $\bregprox$ be closed, convex, and proper with $\operatorname{dom}\bregprox\cap C\neq\emptyset$, and let $\bregsm$ be convex on its domain containing $C$, and differentiable on $C$.
Suppose that $\bregaux=h-\frac{1}{L}\bregsm$ is convex on $C$ for some $L>0$.
Let $x_0 \in C$ and assume that the subproblems in \eqref{eq:bpgm_main} admit selected minimizers in $C$ satisfying the optimality condition \eqref{eq:bpgm_optimality_main}.
Let $x_k$ be the resulting iterates, and let $x_{\compidx}\in\operatorname{dom}F\cap\operatorname{dom}h$ be any reference point satisfying $D_h(x_{\compidx},x_0)<\infty$.
Then, the sequence $\widetilde V_k$ defined in \eqref{eqn:bpgm-lyapunov} satisfies \eqref{eq:bpgm_difference_identity} for $k=0,\dots,N-1$.
Consequently, $\widetilde V_N\le\cdots\le \widetilde V_0=\frac{L}{N}D_h(x_{\compidx},x_0)$, and this implies
\[
    F(x_N)-F(x_{\compidx})\le \frac{L}{N}D_h(x_{\compidx},x_0).
\]
\end{theorem}

\begin{corollary}[Bregman proximal point method]
\label{corollary:bppm-rate}
Let $\bregprox,h$ satisfy the conditions of \cref{theorem:bpgm-lyapunov}.
Let $x_0 \in C$ and let $x_k$ be generated by the Bregman proximal point method
\begin{align}
\tag{BPPM}\label{eq:bppm_main}
    x_{k+1}\in\argmin_{y\in C}
    \left\{
        \bregprox(y)+L D_h(y,x_k)
    \right\}
\end{align}
for $k=0,\dots,N-1$.
Let $x_{\compidx} \in \operatorname{dom}\bregprox\cap\operatorname{dom}h$ satisfy $D_h(x_{\compidx},x_0)<\infty$.
Then the Lyapunov sequence
\[
    \widetilde V_k=\frac{k}{N}\pr{\bregprox(x_k)-\bregprox(x_{\compidx})}+\frac{L}{N}D_h(x_{\compidx},x_k)
\]
is nonincreasing for $k=0,\dots,N$, and consequently,
\[
    \bregprox(x_N)-\bregprox(x_{\compidx})\le \frac{L}{N}D_h(x_{\compidx},x_0) .
\]
\end{corollary}

\begin{lemma}[Tightness of the BPPM rate]
\label{lemma:bppm-tightness}
For any $0<\varepsilon<1$, $N\ge1$, $L>0$, and dimension $d\ge1$, there exist a proper closed convex function $\bregprox:\reals^d\to\reals$, a Legendre-type function $h:\reals^d\to\reals$, and an initial point $x_0\in\reals^d$ such that the iterates of \ref{eq:bppm_main} are uniquely defined and satisfy
\[
    \bregprox(x_N)-\bregprox(x_\star)
    \ge (1-\varepsilon)\frac{L}{N}D_h(x_\star,x_0) 
\]
where $x_\star$ is the unique minimizer of $\bregprox$.
Therefore, the rate of BPPM from \cref{corollary:bppm-rate} 
cannot be improved further.
\end{lemma}

\begin{proof}[Proof of \Cref{lemma:bppm-tightness}]
We first construct a one-dimensional instance.
Fix a parameter $0<\delta<\frac12$, to be chosen sufficiently small at the end.
Define the smoothed absolute value
\[
    s_\delta(x)
    :=
    \begin{cases}
        \dfrac{x^2}{2\delta}+\dfrac{\delta}{2}, & |x|\le\delta,\\[0.6em]
        |x|, & |x|>\delta.
    \end{cases}
\]
Define $\widetilde h_\delta:\reals\to\reals$ by
\[
\widetilde h_\delta(x)
:=
\begin{cases}
\dfrac{N}{2}s_\delta(x+1)-\dfrac{N+2}{2}x-\dfrac{N}{2},
    & x<-\dfrac12,\\[0.8em]
s_\delta(x), & |x|\le\dfrac12,\\[0.8em]
\dfrac{N}{2}s_\delta(x-1)+\dfrac{N+2}{2}x-\dfrac{N}{2},
    & x>\dfrac12.
\end{cases}
\]
A direct differentiation gives
\begin{equation}
\widetilde h_\delta'(x)
=
\begin{cases}
-(N+1), & x\le -1-\delta,\\
-\dfrac{N+2}{2}+\dfrac{N}{2\delta}(x+1),
    & -1-\delta\le x\le -1+\delta,\\[0.4em]
-1, & -1+\delta\le x\le -\delta,\\
\dfrac{x}{\delta}, & |x|\le\delta,\\[0.4em]
1, & \delta\le x\le 1-\delta,\\
\dfrac{N+2}{2}+\dfrac{N}{2\delta}(x-1),
    & 1-\delta\le x\le 1+\delta,\\[0.4em]
N+1, & x\ge 1+\delta.
\end{cases}
    \label{eq:bppm-tight-h-tilde-derivative}
\end{equation}
Thus $\widetilde h_\delta$ is continuously differentiable and convex.
Set
\begin{equation}
    h_\delta(x)
    =
    \frac{1}{L}\left(\widetilde h_\delta(x)+\frac{\delta}{2}x^2\right).
    \label{eq:bppm-tight-h-def}
\end{equation}
The function $h_\delta$ is $\delta/L$-strongly convex, its derivative is strictly increasing, and $h_\delta'(x)\to\pm\infty$ as $x\to\pm\infty$.
Hence $h_\delta$ is a Legendre function, and every BPPM subproblem has a unique minimizer.

Now choose the function and initial point as
\[
    \bregprox(x)=|x|,
    \qquad
    x_0=-1-\delta.
\]
Then $x_\star=0$ is a minimizer of $\bregprox$.
We first show by induction that, for $0\le k\le N$,
\begin{equation}
    x_k\in[-1-\delta,-1+\delta],
    \qquad
    Lh_\delta'(x_k)=-(N+1-k)-\delta(1+\delta).
    \label{eq:bppm-tight-induction}
\end{equation}
The claim is true for $k=0$ by the choice $x_0=-1-\delta$ and \eqref{eq:bppm-tight-h-tilde-derivative}.
Suppose it holds for some $k<N$.
The optimality condition of the BPPM update \eqref{eq:bppm_main} gives
\begin{equation*}
    0\in \partial\bregprox(x_{k+1})
    +L\bigl(h_\delta'(x_{k+1})-h_\delta'(x_k)\bigr),
\end{equation*}
or equivalently
\begin{equation}
    Lh_\delta'(x_{k+1})
    \in
    -\partial\bregprox(x_{k+1})+Lh_\delta'(x_k).
    \label{eq:bppm-tight-optimality}
\end{equation}
Since $-\partial\bregprox(x_{k+1})\subset[-1,1]$, the induction hypothesis implies
\[
    Lh_\delta'(x_{k+1})
    \le
    1+Lh_\delta'(x_k)
    =-(N-k)-\delta(1+\delta)
    <0.
\]
As $Lh_\delta'$ is strictly increasing and $Lh_\delta'(0)=0$, this shows that $x_{k+1}<0$.
Therefore $\partial\bregprox(x_{k+1})=\{-1\}$, and \eqref{eq:bppm-tight-optimality} reduces to
\[
    Lh_\delta'(x_{k+1})
    =
    1+Lh_\delta'(x_k)
    =
    -(N-k)-\delta(1+\delta).
\]
Moreover,
\[
    Lh_\delta'(-1-\delta)=-(N+1)-\delta(1+\delta),
    \qquad
    Lh_\delta'(-1+\delta)=-1-\delta(1-\delta).
\]
Since $0\le k\le N-1$,
\[
    Lh_\delta'(-1-\delta)
    <
    -(N-k)-\delta(1+\delta)
    \le
    Lh_\delta'(-1+\delta).
\]
As $Lh_\delta'$ is increasing, this implies $x_{k+1}\in[-1-\delta,-1+\delta]$.
Together with the identity for $Lh_\delta'(x_{k+1})$ above, this proves \eqref{eq:bppm-tight-induction} with $k$ replaced by $k+1$ and completes the induction.

On the interval $[-1-\delta,-1+\delta]$, \eqref{eq:bppm-tight-h-def} and \eqref{eq:bppm-tight-h-tilde-derivative} give
\begin{equation}
    Lh_\delta'(x)
    =
    -\frac{N+2}{2}+\frac{N}{2\delta}(x+1)+\delta x.
    \label{eq:bppm-tight-h-affine-branch}
\end{equation}
Substituting $x=x_k$ in \eqref{eq:bppm-tight-h-affine-branch} and comparing with \eqref{eq:bppm-tight-induction} gives $\left(\frac{N}{2\delta}+\delta\right)(x_k+1+\delta)=k$, and hence
\begin{equation}
    x_k=-1-\delta+\frac{2\delta}{N+2\delta^2}k,
    \qquad 0\le k\le N.
    \label{eq:bppm-tight-trajectory}
\end{equation}

Because $x_N<0$,
\begin{equation}
    \bregprox(x_N)-\bregprox(x_\star)
    =
    1+\delta-\frac{2\delta N}{N+2\delta^2}.
    \label{eq:bppm-tight-gap}
\end{equation}
For the initial Bregman distance, note that $\widetilde h_\delta(0)=\frac{\delta}{2}$, $\widetilde h_\delta(x_0)=1+(N+1)\delta$, and $\widetilde h_\delta'(x_0)=-(N+1)$, hence $D_{\widetilde h_\delta}(0,x_0)=N+\frac{\delta}{2}$.
The quadratic term in \eqref{eq:bppm-tight-h-def} yields
\begin{equation}
\begin{aligned}
    D_{h_\delta}(x_\star,x_0)
    &=
    \frac{1}{L}\left(
        N+\frac{\delta}{2}
        +\frac{\delta}{2}(1+\delta)^2
    \right)
    =
    \frac{1}{L}\left(
        N+\delta+\delta^2+\frac{\delta^3}{2}
    \right).
\end{aligned}
\label{eq:bppm-tight-divergence}
\end{equation}
Combining \eqref{eq:bppm-tight-gap} and \eqref{eq:bppm-tight-divergence},
\[
\frac{N\bigl(\bregprox(x_N)-\bregprox(x_\star)\bigr)}
     {L D_{h_\delta}(x_\star,x_0)}
=
\frac{N\left(1+\delta-\dfrac{2\delta N}{N+2\delta^2}\right)}
     {N+\delta+\delta^2+\dfrac{\delta^3}{2}}
\longrightarrow 1
\qquad\text{as }\delta\downarrow0.
\]
Thus, for the prescribed $\varepsilon$, one can choose $0<\delta<\frac12$ sufficiently small so that
\[
    \bregprox(x_N)-\bregprox(x_\star)
    \ge
    (1-\varepsilon)\frac{L}{N}D_{h_\delta}(x_\star,x_0).
\]
This proves the claim for dimension $d=1$.

For an arbitrary dimension $d>1$, define, for $z=(z_1,\ldots,z_d)$,
\[
    \widehat{\bregprox}(z)=\bregprox(z_1)+\frac12\sum_{j=2}^d z_j^2,
    \qquad
    \widehat h(z)=h_\delta(z_1)+\frac{1}{2L}\sum_{j=2}^d z_j^2,
\]
and take $z_0=(x_0,0,\ldots,0)$.
Then $\widehat{\bregprox}$ is proper, closed, and convex, $\widehat h$ is Legendre, the unique minimizer is $z_\star=0$, and the BPPM iterates satisfy $z_k=(x_k,0,\ldots,0)$.
Moreover, $D_{\widehat h}(z_\star,z_0)=D_{h_\delta}(x_\star,x_0)$, so the same lower bound holds in dimension $d$.
\end{proof}

\end{document}
